\chapter{Hintergrund}
\label{ch:Hintergrund}

\section{Zeitreihen und Anomalien}

Eine Zeitreihe ordnet Messwerte chronologisch. Bei Passantenfrequenzen hängen
benachbarte Stunden zusammen; zusätzlich wiederholen sich Tagesmuster. Eine
Anomalie ist deshalb nicht zwingend ein großer absoluter Wert, sondern eine
Abweichung vom zeit- und ortsabhängig erwarteten Verhalten. Punktanomalien
betreffen einzelne Stunden. Sprunghafte Niveauwechsel (Level-Shifts),
schleichende Verschiebungen (Drifts) und festgefrorene Sensorwerte (Stuck-at)
erstrecken sich dagegen über mehrere Stunden.

Im vorliegenden Datensatz fehlen vollständige reale Anomaliekennzeichnungen.
Das Lernen des Normalverhaltens erfolgt daher selbstüberwacht: Aus vorhandenen
Messwerten wird ein künstliches Lernziel erzeugt, ohne Klassen wie ``normal'' oder
``anomal'' vorzugeben. Erst der Fehler zwischen Modellschätzung und Messung wird
als Anomaliesignal verwendet.

\section{Erkennung mit Vorhersagen}

Ein Vorhersagemodell erhält die letzten $w$ Zeitpunkte und schätzt den nächsten
Wert $\hat{x}_t$. Ein großer Prognosefehler $|x_t-\hat{x}_t|$ deutet darauf hin,
dass die Beobachtung nicht zum erlernten Verlauf passt. Zusätzliche Messgrößen
können beim Lernen als Kontext dienen, ohne selbst Teil der Alarmentscheidung
zu sein. Vorhersagegüte und Erkennungsleistung sind
dabei nicht gleichbedeutend: Ein kleiner mittlerer Fehler garantiert noch keine klare Trennung seltener
Anomalien.

\section{LSTM, Autoencoder und Transformer}

Long Short-Term Memory (LSTM) ist eine rekurrente Netzstruktur mit steuerbaren
Speicherzuständen und eignet sich zur Verarbeitung geordneter Sequenzen
\cite{hochreiter1997long}. Ein LSTM kann ein Zeitfenster zu einem
Zustandsvektor verdichten und daraus eine Prognose erzeugen.

Ein Autoencoder lernt dagegen, sein Eingabefenster zu rekonstruieren. Ein hoher
Rekonstruktionsfehler kann eine unbekannte Struktur anzeigen. Prognose und
Rekonstruktion beruhen somit auf unterschiedlichen Annahmen über normales
Verhalten. Ein Transformer verarbeitet verschiedene Zeitpunkte mit einem
Aufmerksamkeitsmechanismus und kann dadurch auch weiter auseinanderliegende
Zusammenhänge direkt gewichten. Diese zusätzliche Freiheit ist vor allem bei
längeren Reihen und einer größeren Datenbasis nützlich.
